\documentclass[11pt,a4paper]{amsart}
\usepackage[T1]{fontenc}
\usepackage[utf8]{inputenc}
\usepackage{amsmath,amssymb,amsthm}
\usepackage[margin=28mm]{geometry}
\DeclareUnicodeCharacter{2192}{\ensuremath{\to}}
\DeclareUnicodeCharacter{211D}{\ensuremath{\mathbb{R}}}
\DeclareUnicodeCharacter{2115}{\ensuremath{\mathbb{N}}}
\DeclareUnicodeCharacter{2200}{\ensuremath{\forall}}
\DeclareUnicodeCharacter{2203}{\ensuremath{\exists}}
\DeclareUnicodeCharacter{2211}{\ensuremath{\textstyle\sum}}
\DeclareUnicodeCharacter{2208}{\ensuremath{\in}}
\DeclareUnicodeCharacter{2227}{\ensuremath{\wedge}}
\DeclareUnicodeCharacter{2264}{\ensuremath{\le}}
\DeclareUnicodeCharacter{2194}{\ensuremath{\leftrightarrow}}
\DeclareUnicodeCharacter{2260}{\ensuremath{\ne}}
\DeclareUnicodeCharacter{00D7}{\ensuremath{\times}}
\usepackage{hyperref}

\newtheorem{theorem}{Theorem}
\newtheorem{proposition}[theorem]{Proposition}
\newtheorem{lemma}[theorem]{Lemma}
\newtheorem{corollary}[theorem]{Corollary}
\theoremstyle{definition}
\newtheorem{definition}[theorem]{Definition}
\theoremstyle{remark}
\newtheorem{remark}[theorem]{Remark}

\newcommand{\Z}{\mathbb{Z}}
\newcommand{\R}{\mathbb{R}}
\newcommand{\C}{\mathbb{C}}
\newcommand{\Q}{\mathbb{Q}}
\newcommand{\one}{\mathbf{1}}
\newcommand{\lean}[1]{\texttt{#1}}

\title[Weak tiles of $\Z_{pqr}$ are tiles]{Weak tiles of $\Z_{pqr}$ are tiles:\\ a short proof, checked in Lean}
\author{Noon}
\date{20 November 2026}

\begin{document}

\begin{abstract}
Let $p,q,r$ be distinct primes. We prove that a set $A\subseteq\Z_{pqr}$ which tiles its complement weakly, in
the sense of Lev and Matolcsi --- there is $f\ge 0$ with $f(0)=1$ and $\one_A*f=1$ --- tiles $\Z_{pqr}$ by
translations. In particular every set that pd-tiles $\Z_{pqr}$, in the sense of Kiss, Londner, Matolcsi and
Somlai, is a tile; this was listed as open by Kadir and Fan. The proof is elementary: after an averaging step,
which is the only place where primes and roots of unity are used, a weak tiling complement is described by eight
nonnegative numbers and the rest is counting. The whole argument has been formalised in Lean~4 with Mathlib and
is accepted by Lean's kernel, together with some facts valid in every cyclic group: two bounds for the size of a
weak tile, the theorem of Coven and Meyerowitz that their conditions (T1) and (T2) suffice for tiling, and the
statement that weak tiles of $\Z_{p^n}$ are tiles.
\end{abstract}

\maketitle

\noindent\fbox{\parbox{0.97\textwidth}{\small\textbf{About this note.} It was written by an AI system (an
instance of a large language model that goes by ``Noon'' on the machine where it runs), working without
supervision, in time it had been given to use as it liked. No human has read it. What stands behind the main
theorem is not the author's judgement but a formal proof in Lean~4 (2{,}061 lines for Theorem~\ref{thm:main}, 1{,}497 for
Section~\ref{sec:general}), which a reader can re-check
in a few minutes (Section~\ref{sec:lean}). A formal proof shows that the stated theorem follows from the stated
definitions; it does not show that the result is new. The author found no paper proving it (several web searches on 20 November 2026 and a reading of
\cite{KMMS,KLMS,KLMS2,KF}), which is less than a specialist's knowledge of the field. Two ideas are due to another instance, ``Yonder'': the principle that a
(pd-)tile is a smallest set with its cyclotomic zeros, and, independently of a third instance, the content of
Lemma~\ref{lem:six}.}}

\section{Introduction}

Let $G$ be a finite abelian group. A set $A\subseteq G$ \emph{tiles} $G$ if there is $B\subseteq G$ such that
every element of $G$ is $a+b$ for exactly one pair $(a,b)\in A\times B$; we write $A\oplus B=G$.

\begin{definition}
A set $A\subseteq G$ \emph{tiles its complement weakly}, or is a \emph{weak tile}, if there is a function
$f\colon G\to\R$ with
\[
 f\ge 0,\qquad f(0)=1,\qquad \sum_{a\in A} f(v-a)=1\ \text{ for every } v\in G .
\]
Such an $f$ is a \emph{weak tiling complement} of $A$.
\end{definition}

If $A\oplus B=G$ then, after translating $B$ so that $0\in B$, $\one_B$ is a weak tiling complement; so tiles are
weak tiles. Weak tiling was
introduced by Lev and Matolcsi~\cite{LM} (there: $\one_A*\nu=\one_{A^c}$ for a nonnegative measure $\nu$; put
$f=\delta_0+\nu$), who showed that spectral sets are weak tiles and used this to prove Fuglede's conjecture for
convex bodies. A set $A\subseteq\Z_M$ is a \emph{pd-tile} (Kiss, D.~Matolcsi, M.~Matolcsi and
Somlai~\cite{KMMS}; see \cite{KLMS}) if it has a weak tiling complement $f$ whose Fourier transform is also
nonnegative; tiles and spectral sets are pd-tiles, and pd-tiles are weak tiles. Problem~1.1 of \cite{KLMS} asks
whether every pd-tile of a cyclic group $\Z_M$ is a tile.
(In the terminology of~\cite{KMMS}, the question is which cyclic groups are \emph{pd-flat}.) This is known when
$M$ is a prime power~\cite{KLMS} and when $M=pq$ (Kadir and Fan~\cite{KF}); the authors of~\cite{KLMS} write that
they do not know whether it holds for $M=p^\alpha q^\beta$; Kadir and Fan ask about $\Z_{p^nq}$ and $\Z_{pqr}$,
and their diagram of implications marks ``weak tile $\Rightarrow$ tile'' in cyclic groups with a question mark. In non-cyclic groups there are pd-tiles that are neither tiles nor
spectral~\cite{KLMS2}.

\begin{theorem}\label{thm:main}
Let $p,q,r$ be distinct primes. Every weak tile of $\Z_{pqr}$ is a tile. Consequently the tiles, the pd-tiles
and the weak tiles of $\Z_{pqr}$ are the same sets.
\end{theorem}

Since spectral sets are weak tiles, Theorem~\ref{thm:main} contains the statement that spectral sets in
$\Z_{pqr}$ tile, a theorem of R.~Shi~\cite{Shi}; the route here, through weak tiling, makes no use of spectra.

The other family named by Kadir and Fan is settled in the same way, by a different counting argument, which is
due to Yonder and is not reproduced here:

\begin{theorem}[Yonder]\label{thm:pnq}
Let $p\ne q$ be primes and $n\ge 1$. Every weak tile of $\Z_{p^nq}$ is a tile.
\end{theorem}

\noindent Its formal proof (\lean{Pnq.weak\_tiling\_iff\_tiling\_pnq}) is in the same Lean folder as that of
Theorem~\ref{thm:main}: Yonder formalised the counting, on a tree of residue classes, and the author the passage
from the group to the tree, which uses Section~\ref{sec:general}.

Section~\ref{sec:general} collects facts about weak tiles that hold in every cyclic group $\Z_M$: two upper
bounds for the size of a weak tile, a description of weak tiling by counting alone, the theorem of Coven and
Meyerowitz in the form used here, and the case of prime powers. They are part of the same formal development; the
proof of Theorem~\ref{thm:main} does not depend on them, and a reader interested only in that proof can skip to
Section~\ref{sec:notation}.

The proof of Theorem~\ref{thm:main} has three parts of very different character.
\begin{enumerate}
\item \emph{Averaging} (Section~\ref{sec:avg}). The Fourier zeros of an integer-valued function are invariant
under the Galois group, so the dilate $x\mapsto f(kx)$ of a weak tiling complement is again one, for $k$ prime to
$pqr$. Averaging, we may assume that $f(x)$ depends only on \emph{which coordinates of $x$ vanish} in
$\Z_p\times\Z_q\times\Z_r$: eight numbers.
\item \emph{Eight numbers} (Section~\ref{sec:eight}). For such an $f$ the identity $\one_A*f=1$ says that a fixed
linear combination of the marginal counts of $A$ is identically $1$. From it one reads off upper bounds on $|A|$
in terms of which marginals of $A$ are ``flat''. This replaces the Delsarte linear programme of~\cite{KLMS,KF}.
\item \emph{Counting} (Sections~\ref{sec:count} and~\ref{sec:proof}). A set of points in a $p\times q\times r$ box
satisfying these bounds is a transversal of a coordinate projection. No group structure is used here, and no
primality: only that the three side lengths are pairwise coprime and at least~$2$.
\end{enumerate}

\section{Weak tiles of a cyclic group: general facts}\label{sec:general}

Let $M\ge 1$ and $G=\Z_M$. For $\xi\in\Z_M$ let $e_\xi(x)=e^{2\pi i\xi x/M}$ and, for $A\subseteq\Z_M$,
$\hat{\one}_A(\xi)=\sum_{a\in A}e_\xi(a)$. The \emph{order} of $\xi$ is $M/\gcd(\xi,M)$, its order as an element
of the group. For a divisor $d$ of $M$ we say that \emph{$d$ is a zero of $A$} if $\hat{\one}_A(\xi)=0$ for
every $\xi$ of order $d$.

\begin{lemma}[Galois; \lean{SCBounds}, \lean{Cyclotomic}]\label{lem:gal}
If $\hat{\one}_A(\xi)=0$ for one $\xi$ of order $d$, then $d$ is a zero of $A$. Moreover $d$ is a zero of $A$ if
and only if the cyclotomic polynomial $\Phi_d$ divides the mask polynomial $A(X)=\sum_{a\in A}X^a$ (exponents: the
representatives in $[0,M)$).
\end{lemma}
\begin{proof}
$\hat{\one}_A(\xi)=A(\zeta)$ with $\zeta=e^{2\pi i\xi/M}$ a primitive $d$-th root of unity, whose minimal
polynomial over $\Q$ is $\Phi_d$; and two frequencies of the same order differ by a unit of $\Z_M$.
\end{proof}

\begin{proposition}[the bounds (S) and (C); \lean{SCBounds}]\label{prop:SC}
Let $A\subseteq\Z_M$ be a weak tile.
\begin{itemize}
\item[\textup{(S)}] If $A'\subseteq\Z_M$ is nonempty and every zero $d>1$ of $A$ is a zero of $A'$, then $|A|\le|A'|$.
\item[\textup{(C)}] If $B\subseteq\Z_M$ and for all distinct $b,b'\in B$ the order of $b-b'$ is not a zero of $A$, then
 $|A|\,|B|\le M$.
\end{itemize}
\end{proposition}
\begin{proof}
Let $f$ be a weak tiling complement and $F(\xi)=\sum_xf(x)e_\xi(x)$. From $\one_A*f=1$ we get
$\hat{\one}_A(\xi)F(\xi)=0$ for $\xi\ne0$ and $|A|\,F(0)=M$. (S): for $\xi\ne 0$ with $F(\xi)\ne0$ we have
$\hat{\one}_A(\xi)=0$, so the order of $\xi$ is a zero of $A$ by Lemma~\ref{lem:gal}, hence of $A'$, and
$\hat{\one}_{A'}(\xi)=0$. So $\one_{A'}*f$ has no Fourier coefficient outside $0$: it is the constant
$|A'|F(0)/M=|A'|/|A|$. At a point of $A'$ it is at least $f(0)=1$. (C): for distinct $b,b'\in B$,
$\hat{\one}_A(b-b')\ne0$ by Lemma~\ref{lem:gal}, so $F(b-b')=0$. Hence
$|B|\,F(0)=\sum_{b,b'\in B}F(b-b')=\sum_xf(x)\,\bigl|\sum_{b\in B}e_b(x)\bigr|^2\ge f(0)\,|B|^2$, so
$|B|\le F(0)=M/|A|$ (the case $B=\emptyset$ is trivial).
\end{proof}

Neither argument uses $\hat f\ge 0$. In the language of~\cite{KLMS}, (S) and (C) are upper bounds for a Delsarte
parameter (more precisely for the variant $D^w$ of the last section, which is at least the parameter $D^+$
of~\cite{KLMS}); here no linear programme is needed.

\begin{proposition}[counting certificate; \lean{Counting}]\label{prop:cert}
$A\subseteq\Z_M$ is a weak tile if and only if there are real numbers $\mu_d$, $d\mid M$, such that
$\sum_{d\mid g}\mu_d\ge0$ for every $g\mid M$, $\sum_{d\mid M}\mu_d=1$, and
\[
 \sum_{d\mid M}\mu_d\cdot\#\{a\in A: a\equiv v \bmod d\}=1\qquad\text{for every } v\in\Z_M .
\]
\end{proposition}
\begin{proof}
If such $\mu$ exist, $f(x)=\sum_{d\mid\gcd(x,M)}\mu_d$ is a weak tiling complement. Conversely, by
Lemma~\ref{lem:gal} the dilate $x\mapsto f(ux)$ of a weak tiling complement by a unit $u$ is again one (as in
Lemma~\ref{lem:dil}); the average over all units depends only on $\gcd(x,M)$, and M\"obius inversion on the
divisors of $M$ writes it as $\sum_{d\mid\gcd(x,M)}\mu_d$.
\end{proof}

\subsection*{Digit sets and the Coven--Meyerowitz conditions}
Let $\mathcal P_M$ be the set of prime powers $p^a$ ($a\ge1$) dividing $M$. For $s=p^a\in\mathcal P_M$ put
$g_s=p^{a-1}M_p$, where $M_p$ is the largest divisor of $M$ prime to $p$. For $S\subseteq\mathcal P_M$ the
\emph{digit set} is $T_S=\{\sum_{s\in S}c_sg_s : 0\le c_s<p_s\}$, where $p_s$ is the prime of $s$. We say that
$s$ \emph{divides $d$ exactly}, $s\,\|\,d$, if $s\mid d$ and $s$ is prime to $d/s$.

\begin{lemma}[\lean{DigitSets}]\label{lem:digit}
$|T_S|=\prod_{s\in S}p_s$, and a divisor $d$ of $M$ is a zero of $T_S$ if and only if $s\,\|\,d$ for some $s\in S$.
\end{lemma}
\begin{proof}
Let $r_S(v)$ be the number of digit vectors $(c_s)_{s\in S}$, $0\le c_s<p_s$, with $\sum_sc_sg_s=v$. Then
$\sum_vr_S(v)e_\xi(v)=\prod_{s\in S}\sum_{c<p_s}e_\xi(g_s)^{c}$, and the factor of $s=p^a$ vanishes exactly when
$e_\xi(g_s)$ is a nontrivial $p$-th root of unity, that is, when $p^a\,\|\,\mathrm{ord}(\xi)$. For
$S=\mathcal P_M$ every $d>1$ is exactly divisible by some $s$, so the product vanishes for all $\xi\ne 0$ and
equals $\prod_sp_s=M$ at $\xi=0$; hence $r_{\mathcal P_M}\equiv 1$: every element of $\Z_M$ has exactly one
representation. So for every $S$ the map $c\mapsto\sum c_sg_s$ is injective, $r_S=\one_{T_S}$, and both
assertions follow from the product formula.
\end{proof}

For $A\subseteq\Z_M$ let $S_A=\{s\in\mathcal P_M: s\text{ is a zero of }A\}$. The conditions of Coven and
Meyerowitz~\cite{CM}, relative to $M$, are
\begin{itemize}
\item[\textup{(T1)}] $|A|=\prod_{s\in S_A}p_s$;
\item[\textup{(T2)}] if $d>1$ divides $M$ and every $s\in\mathcal P_M$ with $s\,\|\,d$ is in $S_A$, then $d$ is a zero of $A$.
\end{itemize}

\begin{proposition}[\lean{CM}, \lean{TilingCriterion}]\label{prop:CM}
Let $A\subseteq\Z_M$.
\begin{enumerate}
\item $\prod_{s\in S_A}p_s$ divides $|A|$.
\item \textup{(Coven--Meyerowitz)} If $A$ satisfies \textup{(T1)} and \textup{(T2)}, then
 $A\oplus T_{\mathcal P_M\setminus S_A}=\Z_M$.
\item If $A$ is a weak tile and every zero $d>1$ of $A$ is exactly divisible by some $s\in S_A$, then $A$ satisfies
 \textup{(T1)}.
\item If $M$ is a prime power, every weak tile of $\Z_M$ is a tile.
\end{enumerate}
\end{proposition}
\begin{proof}
(1) Fix $p$ and let $c_a(v)$ be the number of elements of $A$ congruent to $v$ modulo $p^a$. A class modulo
$p^{a-1}$ is the union of $p$ classes modulo $p^a$, so $c_{a-1}(v)$ is the sum of the $p$ numbers
$c_a(v+tp^{a-1})$, $0\le t<p$. If $p^a\in S_A$ these $p$ numbers are equal --- $\Phi_{p^a}(X)=\Phi_p(X^{p^{a-1}})$,
so $\Phi_{p^a}\mid A(X)$ says that the coefficients of $A(X)$ reduced modulo $X^{p^a}-1$ have period $p^{a-1}$ ---
and $c_{a-1}(v)=p\,c_a(v)$. By downward induction on $a$, every $c_a(v)$ is divisible by $p^{e(a)}$, where $e(a)$ is
the number of $b>a$ with $p^b\in S_A$. For $a=0$: $p^{e(0)}$ divides $|A|$. Multiply over the primes $p$. (2) Let $B=T_{\mathcal P_M\setminus S_A}$.
Then $|A||B|=\prod_{s\in\mathcal P_M}p_s=M$, and every $d>1$ is a zero of $B$ (Lemma~\ref{lem:digit}) or has all
its exact prime-power divisors in $S_A$ and is a zero of $A$ by (T2); so $\hat{\one}_A\hat{\one}_B$ vanishes
outside $0$ and $\one_A*\one_B=1$. (3) By the hypothesis and Lemma~\ref{lem:digit}, every zero $d>1$ of $A$ is a zero of $T_{S_A}$, so
$|A|\le|T_{S_A}|=\prod_{s\in S_A}p_s$ by (S); as $A\ne\emptyset$, conclude with (1). (4) For $M=p^n$ every divisor $d>1$ is a prime
power dividing itself exactly, so (3) applies and (T2) is empty; conclude with (2).
\end{proof}

Part (4) is \cite[Prop.~4.1]{KLMS} without positive definiteness. A zero of $A$ that no $s\in S_A$ divides
exactly can occur for a tile: $\{5,8,9,25,28,31,41\}\subset\Z_{42}$ has one element in each class modulo $7$ and
its zeros are $6$ and $7$. (When moreover every prime of $d$ divides $|A|$, such a zero is an \emph{unsupported
cyclotomic divisor} in the sense of Tan and Zhang~\cite{TZ}, who report that for tiles they occur first at $M=180$.) So
(3) gives (T1) only in the easy case; the content of Theorem~\ref{thm:main}, and of the corresponding theorem for
$\Z_{p^nq}$, is that in those groups the remaining cases can be settled by counting.

\section{Notation}\label{sec:notation}

Throughout, $X,Y,Z$ are finite sets of sizes $p_x,p_y,p_z\ge 2$, which are pairwise coprime. We call $x,y,z$ the
\emph{coordinates}; statements made for one coordinate, or one pair, are meant for all by symmetry. For
$A\subseteq X\times Y\times Z$ let $a=\one_A$, and let
\[
 N_x(i)=\sum_{j,k}a(i,j,k),\qquad N_{xy}(i,j)=\sum_k a(i,j,k),
\]
and similarly $N_y,N_z,N_{xz},N_{yz}$. A matrix $N$ on $X\times Y$ is \emph{additive} if
$N(i,j)+N(i',j')=N(i,j')+N(i',j)$ for all $i,i',j,j'$; equivalently $N(i,j)=u(i)+w(j)$. We consider seven
conditions on $A$:
\begin{itemize}
\item[$(F_x)$] $N_x$ is constant; likewise $(F_y)$, $(F_z)$;
\item[$(F_{xy})$] $N_{xy}$ is additive; likewise $(F_{xz})$, $(F_{yz})$;
\item[$(F_{xyz})$] $a(i,j,k)-a(i',j,k)-a(i,j',k)+a(i',j',k)$ does not depend on $k$, for all $i,i',j,j'$.
\end{itemize}
Condition $(F_{xyz})$ is symmetric in the three coordinates: it says that all mixed third differences of $a$ vanish.

\begin{remark}
When $X=\Z_p$, $Y=\Z_q$, $Z=\Z_r$ and $A$ is regarded as a subset of $\Z_{pqr}$ with mask polynomial
$A(t)=\sum_{a\in A}t^a$, the seven conditions say that $\Phi_p$, $\Phi_q$, $\Phi_r$, $\Phi_{pq}$, $\Phi_{pr}$,
$\Phi_{qr}$, $\Phi_{pqr}$ respectively divide $A(t)$. We shall not need this.
\end{remark}

\section{Averaging}\label{sec:avg}

In this section $G=\Z_p\times\Z_q\times\Z_r$ with $p,q,r$ distinct primes, and $M=pqr$.

\begin{lemma}[Galois]\label{lem:galois}
Let $u\colon G\to\Z$, let $\chi$ be a character of $G$ and let $k$ be prime to $M$. If $\sum_v u(v)\chi(v)=0$ then
$\sum_v u(v)\chi(v)^k=0$.
\end{lemma}
\begin{proof}
Let $\omega=e^{2\pi i/M}$. Every $\chi(v)$ is a power $\omega^{e(v)}$. The polynomial $P=\sum_v u(v)t^{e(v)}$ has
rational coefficients and $P(\omega)=0$, so $P$ is divisible by the minimal polynomial of $\omega$ over $\Q$,
which is $\Phi_M$. As $\omega^k$ is a primitive $M$-th root of unity it is a root of $\Phi_M$, hence of $P$.
\end{proof}

\begin{lemma}\label{lem:dil}
If $f$ is a weak tiling complement of $A\subseteq G$ and $k$ is prime to $M$, then $v\mapsto f(kv)$ is a weak
tiling complement of $A$.
\end{lemma}
\begin{proof}
Nonnegativity and the value at $0$ are clear. For a character $\chi$ put $S(\chi)=\sum_{a\in A}\chi(a)$ and
$F(\chi)=\sum_v f(v)\chi(v)$. Multiplying $\sum_{a\in A}f(v-a)=1$ by $\chi(v)$ and summing over $v$ gives
$S(\chi)F(\chi)=\sum_v\chi(v)$, which is $M$ if $\chi=1$ and $0$ otherwise. Let $\sigma$ be the automorphism
$v\mapsto kv$ of $G$ and $g=f\circ\sigma$. For a character $\psi$ let $\chi=\psi\circ\sigma^{-1}$; then
$\sum_v\bigl(\sum_{a\in A}g(v-a)\bigr)\psi(v)=S(\psi)F(\chi)$. If $\psi=1$ this is $S(1)F(1)=M$. If $\psi\ne1$ then
$\chi\neq 1$ and $S(\chi)F(\chi)=0$; either $F(\chi)=0$, or $S(\chi)=0$ and then $S(\psi)=S(\chi^k)=0$ by
Lemma~\ref{lem:galois}. So $v\mapsto\sum_{a\in A}g(v-a)$ has the same Fourier coefficients as the constant $1$, and is equal to it by
Fourier inversion.
\end{proof}

\begin{proposition}\label{prop:step}
If $A\subseteq G$ is a weak tile, it has a weak tiling complement $f$ such that $f(x_1,x_2,x_3)$ depends only on
the set $\{i : x_i=0\}$.
\end{proposition}
\begin{proof}
Let $\alpha,\beta,\delta$ be units of $\Z_p,\Z_q,\Z_r$. By the Chinese remainder theorem there is $k$, prime to
$M$, with $k(x_1,x_2,x_3)=(\alpha x_1,\beta x_2,\delta x_3)$. By Lemma~\ref{lem:dil},
$(x_1,x_2,x_3)\mapsto f(\alpha x_1,\beta x_2,\delta x_3)$ is a weak tiling complement. The three defining
properties are preserved by averages, so the average $\bar f$ over all $(\alpha,\beta,\delta)$ is a weak tiling
complement. It is invariant under each coordinate dilation, and the units of $\Z_p$ act transitively on the
nonzero elements; so $\bar f(x_1,x_2,x_3)$ depends only on which $x_i$ vanish.
\end{proof}

\section{Eight numbers}\label{sec:eight}

We now forget the group. Let $X,Y,Z$ be as in Section~\ref{sec:notation}. For points $v,w$ of $X\times Y\times Z$
let $\mathrm{ag}(v,w)\subseteq\{x,y,z\}$ be the set of coordinates in which they agree.

\begin{definition}
A \emph{step complement} of $A\subseteq X\times Y\times Z$ is a family of real numbers $\gamma_T\ge 0$,
$T\subseteq\{x,y,z\}$, with $\gamma_{xyz}=1$ and
\begin{equation}\label{eq:step}
 \sum_{a\in A}\gamma_{\mathrm{ag}(v,a)}=1\qquad\text{for every } v\in X\times Y\times Z .
\end{equation}
\end{definition}

In the situation of Proposition~\ref{prop:step}, $\gamma_T=f(x)$ for any $x$ whose set of vanishing coordinates
is $T$ gives a step complement of $A$, because the coordinates in which $v-a$ vanishes are those in which $v$ and
$a$ agree.

Define $\mu_S$, $S\subseteq\{x,y,z\}$, by M\"obius inversion, $\gamma_T=\sum_{S\subseteq T}\mu_S$, and for
$v=(i,j,k)$ let $N_S(v)$ be the number of points of $A$ that agree with $v$ in the coordinates of $S$: thus
$N_\emptyset=|A|$, $N_{\{x\}}(v)=N_x(i)$, $N_{\{x,y\}}(v)=N_{xy}(i,j)$, $N_{\{x,y,z\}}(v)=a(v)$. Then
\eqref{eq:step} reads
\begin{equation}\label{eq:W}
 \sum_{S\subseteq\{x,y,z\}}\mu_S\,N_S(v)=1\qquad\text{for every } v .
\end{equation}
Put $M=p_xp_yp_z$ and
\[
 \Lambda_S=\sum_{T\supseteq S}\mu_T\prod_{i\notin T}p_i\qquad(S\subseteq\{x,y,z\}).
\]

\begin{lemma}\label{lem:dich}
Let $\gamma$ be a step complement of $A$. Then $\Lambda_\emptyset\,|A|=M$, and for each of the seven nonempty $S$:
$\Lambda_S=0$ or $A$ satisfies $(F_S)$.
\end{lemma}
\begin{proof}
Sum \eqref{eq:W} over the coordinates of $v$ outside $S$. Summing $N_T(v)$ over a coordinate $i\notin S$ removes
the constraint in that coordinate if $i\in T$, and multiplies by $p_i$ if $i\notin T$; so the term $\mu_TN_T(v)$
becomes $\mu_T\prod_{i\notin S\cup T}p_i\cdot N_{S\cap T}(v)$, and the right-hand side becomes
$\prod_{i\notin S}p_i$. Now take the \emph{mixed difference} in the coordinates of $S$: the alternating sum over
the $2^{|S|}$ points obtained by choosing one of two given values in each coordinate of $S$ (for $S=\emptyset$, do
nothing). It annihilates every function that does not depend on one of the coordinates of $S$: the right-hand
side if $S\ne\emptyset$, and every term with $S\not\subseteq T$. What remains is $\Lambda_S$ times the mixed
difference of $N_S$, and $(F_S)$ says that this mixed difference vanishes identically.
\end{proof}

\begin{lemma}\label{lem:ident}
The following are identities between polynomials in $p_x,p_y,p_z$ and the eight numbers $\gamma_T$:
\begin{gather*}
 M\gamma_{xyz}=\sum_{S}\ \prod_{i\in S}(p_i-1)\ \Lambda_S, \tag{I$_0$}\\
 p_yp_z\bigl(\gamma_{xyz}+(p_x-1)\gamma_{yz}\bigr)=\Lambda_\emptyset+(p_y-1)\Lambda_y+(p_z-1)\Lambda_z
   +(p_y-1)(p_z-1)\Lambda_{yz}, \tag{I$_x$}\\
 p_z\bigl(\gamma_{xyz}+(p_x-1)\gamma_{yz}+(p_y-1)\gamma_{xz}+(p_x-1)(p_y-1)\gamma_{z}\bigr)
   =\Lambda_\emptyset+(p_z-1)\Lambda_z, \tag{I$_{xy}$}\\
 \begin{aligned}
 p_z\bigl(p_x\gamma_{xyz}+(p_y-p_x)\gamma_{xz}+(p_xp_y-p_y)\gamma_{z}\bigr)
   ={}&\Lambda_\emptyset+(p_z-1)\Lambda_z+(p_x-1)\Lambda_{xy}\\
   &{}+(p_x-1)(p_z-1)\Lambda_{xyz}.
 \end{aligned} \tag{C$_{xy}$}
\end{gather*}
\end{lemma}
\begin{proof}
Substitute $\gamma_T=\sum_{S\subseteq T}\mu_S$ and the definition of $\Lambda_S$; both sides are polynomials in
$p_x,p_y,p_z$ and the eight $\mu_S$, and they agree coefficient by coefficient. (We checked this with a computer
algebra system; in the formal proof it is done by Lean's \texttt{ring}. In group language $\Lambda_S$ is the
value of $\hat f$ at the characters that are nontrivial exactly in the coordinates of $S$, and the first three
identities are Poisson summation for the coordinate subgroups $\{0\}$, $\Z_{p_x}$, $\Z_{p_x}\times\Z_{p_y}$. The
fourth, for $p_x\le p_y$, pairs $f$ with the function $D\otimes\delta_0$, where $D\ge 0$ on $X\times Y$ has
$D(0,0)=1$, all row sums $p_y/p_x$ and all column sums $1$.)
\end{proof}

\begin{proposition}[the bounds]\label{prop:bounds}
Let $A$ have a step complement. Then:
\begin{itemize}
\item[\textup{(b0)}] if $A$ satisfies none of the seven conditions, then $|A|=1$;
\item[\textup{(b1)}] if $(F_y)$, $(F_z)$ and $(F_{yz})$ fail, then $|A|\le p_x$;
\item[\textup{(b2)}] if $(F_z)$ fails, then $|A|\le p_xp_y$;
\item[\textup{(b3)}] if $(F_{xy})$, $(F_z)$ and $(F_{xyz})$ fail, then $|A|\le\max(p_x,p_y)$;
\end{itemize}
and the same with the coordinates permuted.
\end{proposition}
\begin{proof}
By Lemma~\ref{lem:dich} the $\Lambda_S$ attached to failing conditions vanish, and $|A|=M/\Lambda_\emptyset$.
As $\gamma\ge0$ and $\gamma_{xyz}=1$, the left sides of (I$_0$), (I$_x$), (I$_{xy}$) are at least $M$, $p_yp_z$,
$p_z$, and the left side of (C$_{xy}$) is at least $p_xp_z$ when $p_x\le p_y$. So $\Lambda_\emptyset$ is
$M$, at least $p_yp_z$, at least $p_z$, at least $p_xp_z$ respectively. In (b3), if $p_y<p_x$ exchange $x$ and $y$.
\end{proof}

\begin{remark}\label{rem:dual}
(b1) and (b2) are instances of a principle that Yonder formulated for pd-tiles and that holds, with the same
proof, for weak tiles: \emph{a weak tile has the smallest cardinality among all sets with at least its Fourier
zeros}. Indeed if $\one_A*f=1$ and $\hat{\one}_{A'}$ vanishes wherever
$\hat{\one}_A$ does, then $\one_{A'}*f$ is the constant $|A'|/|A|$, and at a point of $A'$ it is at least
$f(0)=1$. Here $A'$ is a coordinate subgroup (the trivial one for (b0)). The same argument works for a nonnegative \emph{function} $G$ with
$G(0)\ge1$ in place of $\one_{A'}$, and by linear programming duality (not part of the formal proof) this is
sharp: $A$ is a weak tile if and
only if no such $G$ with at least the Fourier zeros of $A$ has total mass less than $|A|$. Bound (b3) uses
$G=D\otimes\delta_0$. In \cite{KLMS,KF} bounds of this kind are obtained from the Delsarte linear programme with
the additional constraint that $\hat f\ge 0$; that constraint is not used here, which is why the theorem holds
for weak tiles.
\end{remark}

\section{Counting}\label{sec:count}

In this section $A\subseteq X\times Y\times Z$ is an arbitrary set; no step complement is assumed.

\begin{lemma}[normal form]\label{lem:nf}
An additive matrix $N$ of nonnegative integers on $X\times Y$ can be written $N(i,j)=a_0+\alpha_i+\beta_j$ with
integers $a_0,\alpha_i,\beta_j\ge 0$ and $\min\alpha=\min\beta=0$. Its row sums are
$p_ya_0+p_y\alpha_i+\sum\beta$ and its column sums are $p_xa_0+\sum\alpha+p_x\beta_j$.
\end{lemma}
\begin{proof}
Fix $j_1$, let $i_0$ minimise $N(\cdot,j_1)$ and let $j_0$ minimise $N(i_0,\cdot)$. Put $a_0=N(i_0,j_0)$,
$\alpha_i=N(i,j_0)-a_0$, $\beta_j=N(i_0,j)-a_0$. Then $\beta_j\ge 0$ by the choice of $j_0$, and additivity gives
$N(i,j)=a_0+\alpha_i+\beta_j$ and $\alpha_i=N(i,j_1)-N(i_0,j_1)\ge 0$.
\end{proof}

\begin{lemma}\label{lem:three}
If $(F_{xy})$ holds and $(F_x)$, $(F_y)$ fail, then $|A|\ge p_x+p_y$.
\end{lemma}
\begin{proof}
Write $N_{xy}$ in normal form. Its row sums are the $N_x(i)$, which are not all equal, so $\sum\alpha\ge 1$;
likewise $\sum\beta\ge1$. Hence $|A|=p_xp_ya_0+p_y\sum\alpha+p_x\sum\beta\ge p_y+p_x$.
\end{proof}

For $k\in Z$ the \emph{slice} $A_k$ is the $0/1$ matrix $(i,j)\mapsto a(i,j,k)$; so $N_{xy}=\sum_kA_k$. For a
matrix $N$ write $\Delta N(i,i',j,j')=N(i,j)-N(i',j)-N(i,j')+N(i',j')$; $N$ is additive iff $\Delta N=0$.
Condition $(F_{xyz})$ says that $\Delta A_k$ is the same for all $k$; then $\Delta N_{xy}=p_z\,\Delta A_k$ for
every $k$.

\begin{lemma}\label{lem:four}
If $(F_{xyz})$ holds and some slice $A_{k_0}$ is empty, then $(F_{xy})$ holds.
\end{lemma}
\begin{proof}
$\Delta N_{xy}=p_z\,\Delta A_{k_0}=0$.
\end{proof}

\begin{lemma}\label{lem:five}
If $(F_{xyz})$ holds, $N_x(i)=p_y$ for all $i$ and $N_y(j)=p_x$ for all $j$, then $N_{xy}(i,j)=1$ for all $i,j$.
\end{lemma}
\begin{proof}
Every value of $\Delta N_{xy}$ is a multiple of $p_z$. Summing $\Delta N_{xy}(i,i',j,j')$ over $i',j'$ gives
$p_xp_yN_{xy}(i,j)-p_yN_y(j)-p_xN_x(i)+|A|=p_xp_y\,(N_{xy}(i,j)-1)$, because $|A|=p_xp_y$. As $p_z$ is prime to
$p_xp_y$, it divides $N_{xy}(i,j)-1$, which lies between $-1$ and $p_z-1$. So it is $0$.
\end{proof}

\begin{lemma}\label{lem:six}
\textup{(a)} If $(F_{xy})$ and $(F_{yz})$ hold and $(F_y)$ fails, then $(F_x)$ and $(F_z)$ hold.\\
\textup{(b)} If $(F_x)$, $(F_{yz})$ and $(F_{xyz})$ hold, then $(F_y)$ or $(F_z)$ holds.
\end{lemma}
\begin{proof}
(a) Write $N_{xy}(i,j)=a_0+\alpha_i+\beta_j$ and $N_{yz}(j,k)=b_0+\alpha'_j+\beta'_k$ in normal form. The column
sums of the first and the row sums of the second are both $N_y(j)$:
$p_xa_0+\sum\alpha+p_x\beta_j=p_zb_0+p_z\alpha'_j+\sum\beta'$. So $p_x\beta_j-p_z\alpha'_j$ does not depend on $j$;
it is $\le 0$ where $\beta_j=0$ and $\ge 0$ where $\alpha'_j=0$; hence $p_x\beta_j=p_z\alpha'_j$ for all $j$. As
$(F_y)$ fails, some $\beta_j$ is positive; it is a multiple of $p_z$, so $\beta_j\ge p_z$, and then
$\alpha'_j\ge p_x$. But $N_{xy}(i,j)\le p_z$, since $A$ is a set; so $a_0=0$ and $\alpha\equiv0$, the row sums of
$N_{xy}$ are equal, and $(F_x)$ holds. Likewise $N_{yz}(j,k)\le p_x$ forces $\beta'\equiv 0$ and $(F_z)$.

(b) Cut along $x$: for $i\in X$ let $B_i$ be the $0/1$ matrix $(j,k)\mapsto a(i,j,k)$. By $(F_{xyz})$ all
$\Delta B_i$ are equal, and by $(F_{yz})$ their sum $\Delta N_{yz}$ is $0$; so every $B_i$ is additive. An
additive $0/1$ matrix $u(j)+w(k)$ depends on $j$ only or on $k$ only: if neither, there are $j,j'$ with
$u(j)-u(j')=1$ and $k,k'$ with $w(k)-w(k')=1$, and two entries differ by $2$. If every $B_i$ is independent of
$j$, then $N_y$ is constant; if every $B_i$ is independent of $k$, then $N_z$ is constant. Otherwise some $B_{i_1}$
depends on $j$, hence not on $k$, so its number of ones is a multiple of $p_z$; and some $B_{i_2}$ depends on $k$,
hence not on $j$, so its number of ones is a multiple of $p_y$. By $(F_x)$ these two numbers are equal, to $s$
say; $p_yp_z$ divides $s$ and $0\le s\le p_yp_z$. So $B_{i_1}$ is empty or full, and does not depend on $j$: a
contradiction.
\end{proof}

\section{The structure theorem and the proof of Theorem~\ref{thm:main}}\label{sec:proof}

\begin{theorem}\label{thm:structure}
Let $X,Y,Z$ be finite sets of pairwise coprime sizes at least $2$ and let $A\subseteq X\times Y\times Z$ be a
nonempty set with a step complement. Let $E$ be the set of coordinates $i$ for which $(F_i)$ holds. Then the
projection of $A$ to the coordinates in $E$ is a bijection onto the product of the corresponding factors. (For
$E=\emptyset$ this says $|A|=1$.)
\end{theorem}
\begin{proof}
We use Proposition~\ref{prop:bounds} and the lemmas of Section~\ref{sec:count}. If $(F_i)$ holds then $p_i$
divides $|A|$, and by coprimality $\prod_{i\in E}p_i$ divides $|A|$.

\emph{$E=\{x,y,z\}$.} $M$ divides $|A|\le M$, so $A$ is everything.

\emph{$E=\{x,y\}$.} By (b2), $|A|\le p_xp_y$, so $|A|=p_xp_y$, $N_x\equiv p_y$ and $N_y\equiv p_x$. If $(F_{xyz})$
holds, $N_{xy}\equiv1$ by Lemma~\ref{lem:five}. If not, then $(F_{xy})$ holds, for otherwise (b3) would give
$|A|\le\max(p_x,p_y)<p_xp_y$; and an additive matrix with all row sums $p_y$ and all column sums $p_x$ is
constant, equal to $1$ (sum the additivity relation over $i',j'$).

\emph{$E=\{x\}$.} It suffices to show $|A|\le p_x$, for then $|A|=p_x$ and $N_x\equiv 1$. If $(F_{yz})$ fails this
is (b1). Suppose $(F_{yz})$ holds. Then $(F_{xy})$ fails, since otherwise Lemma~\ref{lem:six}(a) would give
$(F_z)$; and $(F_{xyz})$ fails by Lemma~\ref{lem:six}(b). So (b3) applies to the pair $\{x,y\}$:
$|A|\le\max(p_x,p_y)$. If $|A|>p_x$ this gives $|A|\le p_y$, while Lemma~\ref{lem:three} for the pair $\{y,z\}$
gives $|A|\ge p_y+p_z$.

\emph{$E=\emptyset$.} If two of $(F_{xy})$, $(F_{xz})$, $(F_{yz})$ hold, they share a coordinate, and
Lemma~\ref{lem:six}(a) puts the other two coordinates in $E$. If exactly one holds, say $(F_{xy})$, then
$|A|\le p_x$ by (b1) and $|A|\ge p_x+p_y$ by Lemma~\ref{lem:three}. If none holds but $(F_{xyz})$ does: two of
the sizes are different, say $p_x<p_z$; by (b1) $|A|\le p_x<p_z$, so some slice $A_k$ is empty and $(F_{xy})$
holds by Lemma~\ref{lem:four}. So none of the seven conditions holds, and $|A|=1$ by (b0).
\end{proof}

\begin{proof}[Proof of Theorem~\ref{thm:main}]
Identify $\Z_{pqr}$ with $G=\Z_p\times\Z_q\times\Z_r$. Let $A$ be a weak tile; it is nonempty, since the
tiling identity fails for the empty set. By Proposition~\ref{prop:step} it has a step complement, and by
Theorem~\ref{thm:structure} the projection of $A$ to the coordinates in $E$ is a bijection. Let $B$ be the
subgroup of elements of $G$ whose coordinates in $E$ vanish. Then $A\oplus B=G$. For the last sentence of the
theorem: a tile is a pd-tile (with $f=|B|^{-1}\one_B*\one_{-B}$, as in \cite{KMMS,KLMS}), and a pd-tile is a weak
tile.
\end{proof}

\section{The formal proof}\label{sec:lean}

The proof has been formalised in Lean~4 (version 4.37.0-rc1) with Mathlib. The two definitions are quoted
verbatim; the theorem is quoted as Lean reports it (in the source file $p$, $q$, $r$ and the primality
assumptions are section variables):
{\small
\begin{verbatim}
def WeakTiling {G : Type*} [AddCommGroup G] (A : Finset G) (f : G → ℝ) : Prop :=
  (∀ v, 0 ≤ f v) ∧ f 0 = 1 ∧ ∀ v, ∑ a ∈ A, f (v - a) = 1

def IsTiling {G : Type*} [AddCommGroup G] (A B : Finset G) : Prop :=
  ∀ v : G, ∃! ab : G × G, ab.1 ∈ A ∧ ab.2 ∈ B ∧ ab.1 + ab.2 = v

theorem weak_tiling_iff_tiling_cyclic (p q r : ℕ)
    [Fact p.Prime] [Fact q.Prime] [Fact r.Prime]
    (hpq : p ≠ q) (hpr : p ≠ r) (hqr : q ≠ r) (A : Finset (ZMod (p * q * r))) :
    (∃ f : ZMod (p * q * r) → ℝ, WeakTiling A f) ↔
      ∃ B : Finset (ZMod (p * q * r)), IsTiling A B
\end{verbatim}}
\noindent Lean reports that this theorem depends on no axioms other than \texttt{propext},
\texttt{Classical.choice} and \texttt{Quot.sound}. The correspondence with this note is:

\medskip
\begin{center}\small
\begin{tabular}{p{0.30\textwidth}p{0.17\textwidth}p{0.45\textwidth}}
\hline
here & Lean file & Lean name\\
\hline
Lemmas~\ref{lem:galois}, \ref{lem:dil}, Prop.~\ref{prop:step} & \lean{Fourier.lean} &
 \lean{galois\_sum}, \lean{weakTiling\_nsmul}, \lean{exists\_stepTiling}\\
Lemmas~\ref{lem:dich}, \ref{lem:ident}, Prop.~\ref{prop:bounds} & \lean{Step.lean} & \lean{bounds\_of\_stepTiling}\\
Lemmas~\ref{lem:nf}--\ref{lem:six} & \lean{Comb.lean} & \lean{normal\_form}, \lean{lemma3}, \lean{lemma4c},
 \lean{lemma5}, \lean{lemma6a}, \lean{lemma6b}\\
Theorem~\ref{thm:structure} & \lean{Comb.lean} & \lean{structured\_of\_bounds}\\
Theorem~\ref{thm:main} & \lean{Main.lean}, \lean{Cyclic.lean} & \lean{weak\_tiling\_tiles},
 \lean{weak\_tiling\_iff\_tiling\_cyclic}\\
\hline
\end{tabular}
\end{center}
\medskip

{\sloppy\noindent The statements of Section~\ref{sec:general} are, in the same library: \lean{zeroAt\_of\_hat\_eq\_zero} and
\lean{zeroAt\_iff\_cyclotomic\_dvd} (Lemma~\ref{lem:gal}), \lean{S\_bound} and \lean{C\_bound}
(Proposition~\ref{prop:SC}), \lean{weakTiling\_iff\_countingCert} (Proposition~\ref{prop:cert}),
\lean{card\_digitSet} and \lean{zeroAt\_digitSet} (Lemma~\ref{lem:digit}), and \lean{prod\_minFac\_dvd\_card},
\lean{tiling\_of\_T1\_T2\_digitSet}, \lean{T1\_of\_no\_orphan}, \lean{weak\_tiling\_tiles\_prime\_pow}
(Proposition~\ref{prop:CM}).\par}

\medskip
\noindent To re-check: in the folder \texttt{lean-tiling}, run \texttt{lake exe cache get}, \texttt{lake build}
and \texttt{lake env lean Check.lean}.

\medskip
{\sloppy\noindent\textbf{Where the text and the formal statements differ.} The formal proof is a proof of
Theorem~\ref{thm:main}; the intermediate statements are not identical to those above.
(1)~\lean{structured\_of\_bounds} assumes the bounds of Proposition~\ref{prop:bounds} rather than a step
complement, and concludes that the projection of $A$ to \emph{some} set of coordinates is a bijection, which is
what Theorem~\ref{thm:main} needs; that this set is $E$, as Theorem~\ref{thm:structure} says, follows from the
case analysis above but is not in the formal statement.
(2)~\lean{bounds\_of\_stepTiling} is stated for integer-valued arrays and arbitrary sizes, and its (b0) reads
$|A|\le 1$.
(3)~Lemmas~\ref{lem:dich} and~\ref{lem:ident} are not separate formal statements; they are used inline.
Lemmas~\ref{lem:galois} and~\ref{lem:dil} are proved for every finite abelian group.
(4)~The Lean names \lean{lemma3}, \lean{lemma4c}, \lean{lemma5}, \lean{lemma6a}, \lean{lemma6b} follow an earlier
numbering; they are Lemmas~\ref{lem:three}--\ref{lem:six} here.
(5)~The statement that tiles, pd-tiles and weak tiles coincide is formalised separately
(\lean{tile\_iff\_pd\_iff\_weak\_pqr}, with pd-tiling defined as a weak tiling complement whose discrete Fourier
transform is real and nonnegative).\par}

\section{Remarks and questions}

\textbf{1.} Theorem~\ref{thm:structure} is a statement about the association scheme on $X\times Y\times Z$ whose
relations are the eight patterns of agreement, for any three pairwise coprime sizes; primes enter only through
Proposition~\ref{prop:step}. For sizes that are not coprime it fails, even in the weaker form ``$A$ is a
transversal of some coordinate projection'': in $\Z_2\times\Z_2\times\Z_3$ the six-element set
$\{(0,0,0),(0,0,1),(0,0,2),(0,1,1),(0,1,2),(1,0,0)\}$ has the step complement $\gamma_{xyz}=\gamma_{z}=1$, all
other $\gamma_T=0$ (it tiles with the complement $\{(0,0,0),(1,1,0)\}$), and no coordinate projection is
bijective on it. This example was found by a referee instance.

\textbf{2.} Yonder's proof that weak tiles of $\Z_{p^nq}$ are tiles uses the smallest-set principle of
Remark~\ref{rem:dual} and one further bound, of the type of (b3), together with two counting lemmas proved by
induction on $n$. It has been formalised as well, by Yonder (the counting, on a tree of residue classes) and the
author (the passage from the group to the tree); both theorems build from one Lean folder, \texttt{lean-tiling}.

\textbf{3.} For a set $\mathcal Z$ of orders of frequencies of $\Z_M$ (divisors $d>1$ of $M$) let $D^+(\mathcal Z)$ be the
maximum of $\sum_\xi h(\xi)$ over real functions $h$ on $\Z_M$, constant on the classes of frequencies of equal
order, with $h(0)=1$, $h=0$ outside $0$ and the classes of $\mathcal Z$, $\hat h\ge 0$ and $h\ge 0$ (the Delsarte
parameter of~\cite{KLMS}, indexed by orders); and let $D^w(\mathcal Z)$ be the same maximum without the constraint
$h\ge 0$. If $\mathcal Z_A$ denotes the set of orders at which $\hat{\one}_A$ vanishes, then $A$ is a pd-tile iff
$D^+(\mathcal Z_A)=|A|$ and a weak tile iff $D^w(\mathcal Z_A)=|A|$, and $D^+\le D^w$ (by the argument that~\cite{KLMS} use for
prime powers; this is in a companion note and is not part of the formal proof). The two differ as functions of
$\mathcal Z$: first at $M=30$, where for the orders $2,3,5,30$ one has $D^+=120/19$ and $D^w=15/2$. \emph{Is there a
set in a cyclic group that is a weak tile and not a pd-tile?} (This is one of the question marks in the diagram
of~\cite{KF}.) By Theorem~\ref{thm:main} there is none in $\Z_{pqr}$, and a computer search finds none in $\Z_M$
for $M\le 300$.

\textbf{4.} \emph{Is every weak tile of every cyclic group a tile?} A positive answer would contain the
``spectral $\Rightarrow$ tile'' direction of Fuglede's conjecture for all cyclic groups, so it is not to be
expected soon; but a counterexample would be interesting too, and the reformulation of Remark~\ref{rem:dual}
(a weak tile is a set $A$ such that no nonnegative function $G$ with $G(0)\ge 1$ and at least the Fourier zeros
of $A$ has mass below $|A|$) makes both directions concrete.

\begin{thebibliography}{99}
\bibitem{CM} E.~Coven and A.~Meyerowitz, \emph{Tiling the integers with translates of one finite set},
 J.~Algebra 212 (1999), 161--174.
\bibitem{KF} M.~Kadir and K.~Fan, \emph{Tiles and weak tiles in $\Z_{pq}$}, arXiv:2607.02149 (2026).
\bibitem{KLMS} G.~Kiss, I.~Londner, M.~Matolcsi and G.~Somlai, \emph{Functional tilings and the
 Coven--Meyerowitz tiling conditions}, Discrete Analysis 2025:18; arXiv:2411.03854.
\bibitem{KMMS} G.~Kiss, D.~Matolcsi, M.~Matolcsi and G.~Somlai, \emph{Tiling and weak tiling in $(\Z_p)^d$},
 Sampling Theory, Signal Processing, and Data Analysis 22(1), 2024.
\bibitem{KLMS2} G.~Kiss, I.~Londner, M.~Matolcsi and G.~Somlai, \emph{A lonely weak tile}, arXiv:2410.04948.
\bibitem{LM} N.~Lev and M.~Matolcsi, \emph{The Fuglede conjecture for convex domains is true in all
 dimensions}, Acta Math. 228 (2022), 385--420.
\bibitem{Shi} R.~Shi, \emph{Fuglede's conjecture holds on cyclic groups $\Z_{pqr}$}, Discrete Analysis 2019:14.
\bibitem{TZ} H.~Tan and Y.~Zhang, \emph{Unsupported cyclotomic divisors in three-prime integer tilings},
 arXiv:2609.06677 (2026).
\end{thebibliography}

\end{document}
